\begin{summarybox}

	\textbf{\textcolor{CSummary}{一句话核心}}

	旋转椭球表面积由椭圆的半轴长短关系决定：

	长轴在旋转轴方向时，用 $\arcsin$ 公式；短轴在旋转轴方向时，用 $\operatorname{artanh}$ 公式；两轴相等时，退化为球面积公式 $4\pi R^{2}$。

	\medskip

	\textbf{\textcolor{CSummary}{使用条件}}

	\begin{itemize}[leftmargin=*, itemsep=0.5em, topsep=0.4em]

		\item \textbf{条件 1（旋转轴）：}
		      椭圆方程为
		      \[
			      \frac{x^{2}}{a^{2}}+\frac{y^{2}}{b^{2}}=1,
			      \qquad a,b>0.
		      \]
		      椭圆绕 $x$ 轴旋转。

		\item \textbf{条件 2（大小判断）：}
		      必须先判断 $a$ 与 $b$ 的大小关系。

		      如果 $a>b$，则旋转轴方向是长轴方向，使用长椭球公式。

		      如果 $a<b$，则旋转轴方向是短轴方向，使用扁椭球公式。

		      如果 $a=b$，则退化为球，使用球面积公式。

	\end{itemize}

	\medskip

	\textbf{\textcolor{CSummary}{公式汇总}}

	当 $a>b$ 时：
	\[
		S
		=
		2\pi b^{2}
		+
		\frac{2\pi ab}{e}\arcsin e,
		\qquad
		e=\sqrt{1-\frac{b^{2}}{a^{2}}}.
	\]

	当 $a<b$ 时：
	\[
		S
		=
		2\pi b^{2}
		+
		\frac{2\pi a^{2}}{e}\operatorname{artanh}e,
		\qquad
		e=\sqrt{1-\frac{a^{2}}{b^{2}}}.
	\]

	其中
	\[
		\operatorname{artanh}e
		=
		\frac12\ln\frac{1+e}{1-e}.
	\]

	当 $a=b=R$ 时：
	\[
		S=4\pi R^{2}.
	\]

	\medskip

	\textbf{\textcolor{CSummary}{关键提醒}}

	\begin{itemize}[leftmargin=*, itemsep=0.5em, topsep=0.4em]

		\item \textbf{易错点（公式混淆）：}
		      不要看到椭圆方程就直接套公式。必须先判断 $a>b$、$a<b$，还是 $a=b$。

		\item \textbf{检查项（离心率定义）：}
		      不同情形下，离心率 $e$ 的表达式不同。

		      当 $a>b$ 时：
		      \[
			      e=\sqrt{1-\frac{b^{2}}{a^{2}}}.
		      \]

		      当 $a<b$ 时：
		      \[
			      e=\sqrt{1-\frac{a^{2}}{b^{2}}}.
		      \]

		      本质上，分母应该是较大的半轴平方。

	\end{itemize}

\end{summarybox}